\documentclass{article}
\usepackage[a4paper, total={6.5in, 10in}]{geometry}
\setlength{\parindent}{0pt}
\usepackage{tcolorbox}
\tcbset{colback=white!94!black, colframe=black!60!white, coltitle = white, boxrule=0.8pt, arc=2mm}
\newtcolorbox{boxs}[2][]{title= #2,#1}
\usepackage{datetime2}
\usepackage[british]{babel}
\usepackage{amsfonts}
\usepackage{amsmath}

\title{\textbf{Workings from Bartle-Sherbet and Lebl}}
\author{Robert O'Connell}
\date{\today}
\begin{document}
\maketitle
\vspace{5mm}

\section*{Lebl}
\subsection*{Proof of Strong Induction}
Let \(P(n)\) be statement depending on a natural number n. Suppose that

\begin{itemize}
    \item \(P(1)\) is true
    \item If \(P(k)\) is true for all \(1,2, \dots k\) then \(P(k+1)\) is true
\end{itemize}

Let \(S\) be the set of natural numbers \(s\) such that \(P(s)\) is not true.
Suppose for contradiction that \(S\) is non-empty, then \(S\) has a least element, say \(m\), by the well ordering property.
We know \(1 \notin S\) since \(P(1)\) is true. Then \(m > 1\) and \(m - 1\) is a natural number as well.
Then \(P(m-1)\) must be true since \(m\),  is the least element for which \(P(s)\) is not true, this also is true for all \(s<m\), but then \(P(m-1+1)\) must be true.
A contradiction, thus m is not in \(S\), meaning there is not a least element in \(S\), i.e. it is empty.
Thus \(P(k)\) is true for all natural numbers
\\
\(\square\)\footnote{meh}

\newpage
\section*{Bartle-Sherbet}
\subsection*{Exercise 1.1.7}
Let \(A := B := {x \in \mathbb{R}: -1 \leq x \leq 1}\) and consider the subset
\(C := (x,y):x^2 +y^2 = 1\) of \(A \times B\). Is this set a function? Explain.
\\

If \(C\), a subset of \(A \times B\), is a function from \(A \rightarrow B\), for
each \(a \in A\), there exists a unique \(b \in B\) with \((a,b) \in C\)

For each \(x \in A\) the function which gives the \(y \in B\) is \(y = \pm \sqrt{1 - x^2}\),
for this \(y\) to not be in \(B\), we either need \(1 - x^2 < 0\) or \(|\sqrt{1-x^2}| > 1\).

The restrictions these put on \(x\) are \(-1 < x < 1 \) which are all in the set \(A\)

But then a problem arise with the uniqueness of the \(y\), when \(x \neq \pm 1\), the function has two possible
values, say \(x = 0.5\) then \(y = \pm \frac{\sqrt{3}}{2}\). A contradiction to our defintion
\\

Thus C is not a function.

\end{document}